\begin{table}[H]
\centering
\caption{Top 10 Departamentos por Defección del Partido Colorado al Frente Amplio (Modelo con PI)}
\label{tab:dept_pc_top10}
\small
\begin{adjustbox}{max width=\textwidth}
\begin{tabular}{lrrrcc}
\toprule
\textbf{Departamento} & \textbf{N} & \textbf{Votos PC} & \textbf{PC$\rightarrow$FA (\%)} & \textbf{IC 95\%} & \textbf{Votos Est.} \\
\midrule
Río Negro & 121 & 8,148 & 50.3 & [47.3, 53.3] & 4,102 \\
Salto & 288 & 18,940 & 20.4 & [16.3, 24.2] & 3,866 \\
Flores & 59 & 4,083 & 19.3 & [6.6, 31.4] & 789 \\
Rivera & 263 & 27,280 & 19.3 & [16.4, 22.0] & 5,266 \\
Durazno & 152 & 9,341 & 16.8 & [11.0, 21.9] & 1,565 \\
Tacuarembó & 234 & 17,012 & 12.2 & [7.8, 16.1] & 2,070 \\
Florida & 179 & 9,125 & 10.5 & [3.0, 17.8] & 954 \\
Canelones & 1,095 & 55,500 & 9.8 & [6.1, 12.6] & 5,412 \\
San José & 252 & 12,643 & 7.1 & [2.2, 11.8] & 892 \\
Lavalleja & 144 & 10,935 & 5.0 & [0.6, 10.1] & 546 \\
\bottomrule
\end{tabular}
\end{adjustbox}
\begin{flushleft}
\small \textit{Nota:} Estimaciones del modelo con Partido Independiente (PI) como categoría separada. Departamentos ordenados por probabilidad de transición PC$\rightarrow$FA. Votos estimados = votos PC $\times$ tasa defección. Todos los modelos convergieron ($\hat{R} < 1.02$).
\end{flushleft}
\end{table}
